% STARTHERE

La struttura seguente ricalca la guideline di reporting GRRAS (i
numeri tra parentesi graffe, che si consiglia di lasciare, indicano
gli item della stessa); ai fini della compilazione delle parti di
propria competenza si consiglia di fare riferimento all'\href{https://linkinghub.elsevier.com/retrieve/pii/S0895-4356(10)00097-1}{articolo relativo}

\section{Title/Abstract \{1\}}

Identify in title or abstract that interrater/intrarater reliability
or agreement was investigated.

\section{Introduction}

\subsection{Diagnostic/measurement devices \{2\}}

Name and describe the diagnostic or measurement device of interest
explicitly.

\subsection{Subject population \{3\}}

Specify the subject population of interest.

\subsection{Rater population \{4\}}

Specify the rater population of interest (if applicable).

\subsection{Background/rationale \{5\}}

Describe what is already known about reliability and agreement and
provide a rationale for the study (if applicable).

\section{Methods}

\subsection{Sample size \{6\}}

Explain how the sample size was chosen. State the determined number of
raters, subjects/objects, and replicate observations.

\subsection{Sampling method \{7\}}

Describe the sampling method.

\subsection{Measurement/rating process \{8\}}

Describe the measurement/rating process (e.g. time interval between
repeated measurements, availability of clinical information,
blinding).

\subsection{Blindness/indipendence \{9\}}

State whether measurements/ratings were conducted independently.

\subsection{Statistical analysis \{10\}}

Describe the statistical analysis.

\section{Results}

\subsection{Observation numerosity \{11\}}

State the actual number of raters and subjects/objects which were
included and the number of replicate observations which were
conducted.

\subsection{Raters and subjects\{12\}}

Describe the sample characteristics of raters and subjects
(e.g. training, experience).

\subsection{Estimates \{13\}}

Report estimates of reliability and agreement including measures of
statistical uncertainty.

\section{Discussion \{14\}}

Discuss the practical relevance of results.

\section{Auxiliary material \{15\}}

Provide detailed results if possible (e.g. online).
